% page 363 du fac-simile (image p-362 du scan)
% bloc 165.1 x 229.0 mm ; marge gauche 36.9 mm
\pgc{15.240mm}{0.000mm}{
\l{\sou{manometro}. Instrumento per qua onu mezuras la tenseso o la}
\l{\cel{3}elastikeso di la gasi, di la vapori. - DEFIRS.}
\l{}
\l{\sou{manoto}. Ligilo kordo o fera, per qua onu interligas la manui}
\l{\cel{3}di homo kaptita por impedar lu eskapar od agar malaji.}
\l{\cel{3}FIS.}
\l{}
\l{\sou{manovrar}. I. (\sou{trans.}) Funcionigar normale, per manuo (instru-}
\l{\cel{3}mento, mashino). - II. (\sou{netran}s.) (I) (\sou{alud}ante\cel{1}\sou{armei, navi})}
\l{\cel{3}movar, evolucionar.(2) (\sou{aludante la persono qua imperas}\cel{1}ad}
\l{\cel{3}\sou{armei)}\cel{2}Evolucionigar armeo, eskadro. - III. (\sou{netrans.})}
\l{\cel{3}Kombinar serio de agi, de demarshi por atingar skopo. - DEFIRS}
\l{}
\l{\sou{mansardo}. I.Tekto di qua la karpenturo esas quaze ruptita. -}
\l{\cel{3}II. Chambro facita sub la quaze-ruptita karpenturo, e di}
\l{\cel{3}qua la parto supra esas plu streta pro la inklineso di la}
\l{\cel{3}tekto. - DFIR.}
\l{}
\l{\sou{mansheto}. Garnituro, sutita o muntita kom ornivo ye la extremajo}
\l{\cel{3}di la maniki dikamizo, e qua salias de la vesto, de la robo,}
\l{\cel{3}tale ke la karpo kovresas. - DFIR.}
\l{}
\l{\sou{mantear}. (\sou{trans.})\cel{2}Molestar (ulu) per igar lua korpo saltar de}
\l{\cel{3}lito-kovrilo quan sukusas quar personi. - S.}
\l{}
\l{\sou{mantelo}.\cel{2}Vesto ampla e longa, sen maniki, quan la homi +weras}
\l{\cel{3}sur lia vesti cetera. - DEFIRS.}
\l{}
\l{mantenar. (\sou{trans.}) I. Tenar ye maniero tala ke nula movo lua}
\l{\cel{3}esas posibla. - II. Tenar en stando tala ke nula chanjo o}
\l{\cel{3}vario esas posibla. - EFIS.}
\l{}
\l{\sou{mantilio}. Peco de stofo, de dentelo nigra, quan la Hispanini}
\l{\cel{3}+weras sur sua kapo e qua retrofalas adsur la shultro.-}
\l{\cel{3}DEFIRS.}
\l{}
\l{mantiso.\cel{2}(\sou{matem.}) La parto decimala di logaritmo.-}
\l{}
\l{manuo. (\sou{anat.)}\cel{2}Organo di prenado e di tushado, ye la extremajo}
\l{\cel{3}di la brakio, che la homo, e kun kin fingri de qui la kin}
\l{\cel{3}esma (la polexo) povas opozar su a la quar ceterá.EFIS.}
\l{}
\l{manufakturo.\cel{2}Granda laboreyo en qua fasonesas, manue, ula pro-}
\l{\cel{3}dukturi qui postulas kelka delikateso laborala. (La manufak-}
\l{\cel{3}turi sencese remplasesas da \textquotedbl{}fabrikeyi\textquotedbl{}, segun ke la mashini}
\l{\cel{3}substitucesas a la manuolaboro). - DEFRS.}
\l{}
\l{\cel{1}\sou{manuskripto}. Verko skribita manue o per skribmashino, ankore}
\l{\cel{3}ne imprimita. - DEFIRS.}
\l{}
\l{\cel{1}\sou{manzanilo}. (\sou{bot.})\cel{2}Arboro di Antili, ek la familio\textquotedbl{}euforbiacei\textquotedbl{},}
\l{\cel{4}di qua la stipo e la frukto kontenas veneno. -L.\cel{1}\sou{hippomano}}
\l{\cel{4}- DEFIS.}
}
